\documentclass[11pt,a4paper,oneside,openright]{book}
\usepackage[spanish,activeacute,es-noshorthands]{babel}
\usepackage[utf8]{inputenc}
\usepackage[T1]{fontenc}
\usepackage{lmodern}
\usepackage{amsmath, amssymb, amsthm}
\usepackage{graphicx}
\usepackage[dvipsnames]{xcolor}
\usepackage{tikz}
\usepackage{listings}
\usepackage{geometry}
\geometry{margin=2.5cm, headheight=14pt}
\usepackage{microtype}
\usepackage{booktabs}
\usepackage{longtable}
\usepackage{enumitem}
\usepackage{fancyhdr}
\usepackage{tcolorbox}
\tcbuselibrary{skins, breakable}
\usepackage{caption}
\captionsetup{font=small, labelfont=bf}
\usepackage{hyperref}
\hypersetup{
  colorlinks=true, linkcolor=NavyBlue, citecolor=ForestGreen, urlcolor=BrickRed,
  pdftitle={M0 --- CubeSat 101 --- introduccion al ecosistema},
  pdfauthor={INTISAT --- CubeSat-EdgeAI-Payload}
}

\newtheorem{definicion}{Definicion}[chapter]
\newtheorem{teorema}[definicion]{Teorema}
\newtheorem*{nota}{Nota}

\definecolor{codebg}{RGB}{248, 249, 250}

\newtcolorbox{intuibox}[1][]{
  colback=green!5, colframe=green!40!black,
  fonttitle=\bfseries, title=Intuicion,
  breakable, sharp corners=south, #1
}
\newtcolorbox{calculo}[1][]{
  colback=gray!4, colframe=gray!55,
  fonttitle=\bfseries, title=Calculo,
  breakable, sharp corners=south, #1
}
\newtcolorbox{payload}[1][]{
  colback=blue!4, colframe=blue!50!black,
  fonttitle=\bfseries, title=Conexion con el payload INTISAT,
  breakable, sharp corners=south, #1
}
\newtcolorbox{cuidado}[1][]{
  colback=red!5, colframe=red!50!black,
  fonttitle=\bfseries, title=Cuidado,
  breakable, sharp corners=south, #1
}
\newtcolorbox{ejercicio}[1][]{
  colback=orange!5, colframe=orange!60!black,
  fonttitle=\bfseries, title=Ejercicios,
  breakable, sharp corners=south, #1
}

\pagestyle{fancy}
\fancyhf{}
\fancyhead[LE,RO]{\thepage}
\fancyhead[RE]{\nouppercase{\leftmark}}
\fancyhead[LO]{\nouppercase{\rightmark}}

\title{\Huge\bfseries M0 --- CubeSat 101\\[0.4em]
\Large Introduccion al ecosistema de los nanosatelites\\[0.6em]
\large Manual de la coleccion CubeSat (M0 / 11)}
\author{Coleccion CubeSat para el proyecto INTISAT}
\date{\today}

\begin{document}

\frontmatter
\maketitle

\chapter*{Prefacio}
\addcontentsline{toc}{chapter}{Prefacio}

Este es el manual numero \textbf{cero} de la coleccion CubeSat. Su
proposito es darte el lenguaje y los conceptos para entender los manuales
que siguen.

Si estas leyendo esto, probablemente sos parte de un equipo que quiere
construir un CubeSat (o ya esta en proceso) y necesitas \emph{ubicarte en
el panorama}. Te lo voy a contar como una historia: que es esto, de donde
salio, como funciona, quien lo usa, y donde encaja tu proyecto.

\section*{Audiencia}

\begin{itemize}[leftmargin=2em]
  \item Estudiantes de ingenieria que no han trabajado en aeroespacial.
  \item Profesionales de otras industrias (mineria, software, electronica)
        que estan transicionando.
  \item Miembros nuevos del equipo INTISAT que vienen sin background.
  \item Curiosos en general.
\end{itemize}

\section*{Prerequisitos}

\textbf{Ninguno tecnico}. Algo de matematica de secundaria ayuda para los
calculos de orbita. Si entendes Pitagoras y la formula de un circulo,
estas bien.

\section*{Como leer}

El manual son \textbf{8 capitulos} mas apendices. Total $\sim$ 50 paginas.
Lectura completa: 1 dia laboral. Lectura selectiva: capitulos sueltos
segun necesidad.

\textbf{Cajas de color} (igual convencion que el resto de la coleccion):
\begin{itemize}[leftmargin=2em]
  \item \textbf{Verde} (Intuicion): analogias que ayudan a entender.
  \item \textbf{Gris} (Calculo): cuentas y derivaciones.
  \item \textbf{Azul} (Conexion con payload): como aplica al INTISAT.
  \item \textbf{Rojo} (Cuidado): errores comunes a evitar.
  \item \textbf{Naranja} (Ejercicios): preguntas para verificar.
\end{itemize}

\tableofcontents

\mainmatter

% =====================================================================
\chapter{Que es un CubeSat}
\label{ch:que-es}

\section{Definicion oficial}

\begin{definicion}[CubeSat]
Un \emph{CubeSat} es un satelite pequeño cuyo diseño cumple con el
estandar \textbf{CubeSat Design Specification (CDS)} desarrollado por
California Polytechnic State University (Cal Poly) en 1999.
\end{definicion}

El estandar define las dimensiones, masa y propiedades mecanicas que
permiten que muchos CubeSats compartan el mismo lanzador en un dispenser
estandar. Ese fue su gran aporte: \textbf{estandarizar el formato} para
abaratar el acceso al espacio.

\section{La unidad: 1U}

La unidad basica es \textbf{1U} (de "unit"):

\begin{itemize}[leftmargin=2em]
  \item Cubo de $10 \times 10 \times 10$ cm.
  \item Masa maxima: $1.33$ kg (originalmente $1.0$ kg, luego revisado).
  \item Centro de gravedad: dentro de un cubo de $\pm 2$ cm desde el centro
        geometrico.
  \item Materiales aprobados: aluminio 6061-T6 o 7075-T73.
  \item Estructura externa con rails de aluminio en las 4 aristas
        longitudinales.
\end{itemize}

\begin{intuibox}
Pensa el 1U como una caja de zapatos chica (de un par de tenis de niño)
pero hecha de aluminio. Entra en la palma de la mano. Pesa lo mismo que
una botella de gaseosa llena.
\end{intuibox}

\section{Multiplos: 1.5U, 2U, 3U, 6U, 12U, 16U, 27U}

Los CubeSats se construyen apilando unidades:

\begin{center}
\begin{tabular}{lll}
\toprule
\textbf{Formato} & \textbf{Dimensiones} & \textbf{Masa max} \\
\midrule
1U     & $10 \times 10 \times 10$ cm     & 1.33 kg \\
1.5U   & $10 \times 10 \times 15$ cm     & 2.0 kg  \\
2U     & $10 \times 10 \times 22.7$ cm   & 2.66 kg \\
3U     & $10 \times 10 \times 34$ cm     & 4.0 kg  \\
6U     & $20 \times 10 \times 34$ cm     & 12.0 kg \\
12U    & $20 \times 20 \times 34$ cm     & 24.0 kg \\
16U    & $20 \times 20 \times 45$ cm     & 32.0 kg \\
27U    & $30 \times 30 \times 34$ cm     & 54.0 kg \\
\bottomrule
\end{tabular}
\end{center}

\textbf{Tipico hoy}: el 80\% de los CubeSats lanzados son 3U o 6U.

\begin{payload}
El INTISAT es probablemente un \textbf{3U} (a confirmar con el equipo de
mecanica). Tu payload de microscopia ocupa un slot de 1U dentro del
3U, dejando 1U para OBC y 1U para EPS + COMMS.
\end{payload}

\section{Por que existen}

Antes de los CubeSats, lanzar un satelite costaba decenas o cientos de
millones de dolares. Solo gobiernos y grandes corporaciones podian.

CubeSat baja la barrera a:
\begin{itemize}[leftmargin=2em]
  \item \textbf{Costo}: USD 100,000 -- 500,000 para una mision academica
        completa (lanzamiento incluido), gracias al rideshare.
  \item \textbf{Tiempo}: 1-3 años de desarrollo (vs 7-10 de satelite tradicional).
  \item \textbf{Riesgo}: tolerable. Una mision fallida no quiebra la
        organizacion.
  \item \textbf{Educacion}: estudiantes universitarios pueden hacer una
        mision completa antes de graduarse.
\end{itemize}

\section{Lo que NO es un CubeSat}

\begin{cuidado}
\textbf{No todo satelite pequeño es CubeSat}. Un nanosatelite que no
cumple CDS no es CubeSat, aunque pese menos de 10 kg. Casos:
\begin{itemize}[leftmargin=1em]
  \item TubeSat (cilindricos, hasta 0.75 kg) --- otro estandar.
  \item PocketQube (5x5x5 cm) --- otro estandar.
  \item Satellogic NewSat-1 (38 kg) --- nanosatelite pero no CubeSat.
\end{itemize}
\end{cuidado}

\section{Numeros del ecosistema (2024)}

\begin{itemize}[leftmargin=2em]
  \item Mas de \textbf{2 000} CubeSats lanzados desde 1999.
  \item Tasa actual: $\sim 200$ por año.
  \item Pais con mas: Estados Unidos ($\sim 60\%$).
  \item Universidades involucradas: $> 100$ en el mundo.
\end{itemize}

% =====================================================================
\chapter{Historia breve}
\label{ch:historia}

\section{El origen --- 1999}

En 1999, Jordi Puig-Suari (Cal Poly) y Bob Twiggs (Stanford) proponen un
estandar para que los estudiantes puedan hacer satelites \emph{de verdad}.
La motivacion era educativa: en aquel entonces, la unica forma de
participar en una mision real era hacer una pasantia en NASA o JPL.

\section{Hitos clave}

\begin{center}
\begin{tabular}{p{1.3cm}p{12cm}}
\toprule
\textbf{Año} & \textbf{Evento} \\
\midrule
1999  & Cal Poly y Stanford definen el estandar CDS. \\
2003  & Primer lanzamiento exitoso (6 CubeSats al mismo tiempo, ruso). \\
2006  & GeneSat-1 (NASA Ames): primer CubeSat con biologia (E. coli). \\
2009  & PharmaSat (NASA Ames): levaduras + microscopia de fluorescencia. \\
2013  & PhoneSat (NASA): smartphone Nexus como OBC. \\
2014  & Planet Labs lanza la primera flota de Doves (constelacion EO). \\
2017  & UPSat (Libre Space): primer CubeSat 100\% open source. \\
2018  & MarCO (NASA): primeros CubeSats interplanetarios (Marte). \\
2020  & PhiSat-1 (ESA): primer satelite con IA on-board. \\
2022  & HYPSO-1 (NTNU): IA on-board con RPi CM4 (parecido al INTISAT). \\
2024  & PhiSat-2 (ESA): 6 apps de IA simultaneas. \\
2024  & CogniSAT-6 (Aiko): primer CubeSat EO comercial con AI integrada. \\
\bottomrule
\end{tabular}
\end{center}

\section{Los antecedentes peruanos}

\begin{itemize}[leftmargin=2em]
  \item \textbf{CHASQUI-I} (UNI, 2014): primer CubeSat peruano, lanzado
        desde la ISS.
  \item \textbf{PUCP-Sat 1} (PUCP, 2013): primer CubeSat operativo del pais.
  \item \textbf{POCKET-PUCP} (PUCP, 2013): el "femtosat" mas chico de su tiempo.
  \item \textbf{UAPSAT-1} (UAP, 2013): educativo.
  \item \textbf{INTISAT} (en curso): tu proyecto.
\end{itemize}

% =====================================================================
\chapter{Anatomia de un CubeSat}
\label{ch:anatomia}

\section{Los 6 subsistemas}

Todo CubeSat tiene los mismos 6 subsistemas, sin importar la mision:

\begin{center}
\begin{tabular}{p{2.5cm}p{12.5cm}}
\toprule
\textbf{Subsistema} & \textbf{Funcion} \\
\midrule
\textbf{Mecanica}   & estructura, aleta, deployments. Es el "chasis". \\
\textbf{EPS}        & Electrical Power System: paneles solares, bateria,
                      reguladores. La energia. \\
\textbf{OBC}        & On-Board Computer: la "cabeza" que coordina todo. \\
\textbf{ADCS}       & Attitude Determination \& Control System: hacia
                      donde apunta el satelite. \\
\textbf{COMMS}      & comunicacion con tierra (UHF, S-Band, etc.). \\
\textbf{Payload}    & la mision cientifica o comercial (lo que justifica
                      lanzarlo). \\
\bottomrule
\end{tabular}
\end{center}

\section{Mecanica}

La estructura es de aluminio anodizado (negro). Cuatro rails
longitudinales sobresalen para encajar en el dispenser. Adentro, slots
PC-104 (estandar de bus electrico/mecanico) sostienen las placas de cada
subsistema.

\textbf{Cosas mecanicas a considerar}:
\begin{itemize}[leftmargin=2em]
  \item Centro de gravedad bien centrado (criterio CalPoly).
  \item Vibracion durante el lanzamiento (test 5--2000 Hz).
  \item Outgassing (materiales que liberen vapor en vacio).
  \item Despliegues post-lanzamiento (antenas, paneles solares).
\end{itemize}

\section{EPS --- el corazon energetico}

\begin{itemize}[leftmargin=2em]
  \item \textbf{Paneles solares}: tipicamente cubren 4 caras del 3U.
        Generan 6--10 W promedio.
  \item \textbf{Bateria}: Li-ion, 10--40 Wh.
  \item \textbf{Reguladores}: generan 3.3V, 5V, 12V para los demas
        subsistemas.
  \item \textbf{Distribucion}: rails con proteccion contra cortos.
\end{itemize}

\section{OBC --- el cerebro}

Generalmente un microcontrolador (STM32, ESP32) o un SoC (Raspberry Pi CM4
en los modernos). Corre el flight software que coordina todo: lee
sensores de actitud, controla actuadores, gestiona comms, dispara payload.

\section{ADCS --- la orientacion}

Tres formas de controlar la actitud:
\begin{itemize}[leftmargin=2em]
  \item \textbf{Detumble} con magnetorquers (electroimanes que interactuan
        con el campo magnetico terrestre).
  \item \textbf{Pointing} fino con reaction wheels (ruedas que aceleran
        para girar el satelite).
  \item \textbf{Determinacion} con gyros, magnetometros, sun sensors, star
        trackers.
\end{itemize}

\section{COMMS --- el enlace}

\begin{itemize}[leftmargin=2em]
  \item \textbf{UHF} (430 MHz): comms basicos, amateur radio. Baja velocidad
        (1--9.6 kbps), pero atraviesa nubes y barato.
  \item \textbf{S-Band} (2.2 GHz): comms medios. 100 kbps -- 1 Mbps.
  \item \textbf{X-Band} (8 GHz): downlink rapido. 10--100 Mbps.
\end{itemize}

\section{Payload --- la mision}

Es lo que el CubeSat \emph{hace}. Puede ser:
\begin{itemize}[leftmargin=2em]
  \item Camara de observacion de la Tierra (Planet, Spire).
  \item Sensor cientifico (campo magnetico, radiacion, gases).
  \item Demostrador tecnologico (nuevo propulsor, nuevo material).
  \item \textbf{Microscopia biologica con IA on-board (INTISAT)}.
\end{itemize}

\begin{payload}
Tu payload del INTISAT es:
\begin{itemize}[leftmargin=1em]
  \item Raspberry Pi 5 como computador del payload.
  \item Sensor OmniVision OV5647 sin lente (lensless).
  \item OLED SSD1351 como matriz de iluminacion programable.
  \item Mascaras opticas intercambiables (pinhole, cuadrada, slit).
  \item Pipeline de IA para segmentacion celular (StarDist, Cellpose).
\end{itemize}
\end{payload}

\section{Cuanto pesa y cuanto consume cada uno}

\textbf{Para un 3U tipico}:

\begin{center}
\begin{tabular}{lcc}
\toprule
\textbf{Subsistema} & \textbf{Masa (g)} & \textbf{Consumo idle (W)} \\
\midrule
Mecanica + harness   & 800  & 0 \\
EPS                  & 400  & 0.2 \\
OBC                  & 100  & 1.5 \\
ADCS                 & 600  & 0.5 \\
COMMS                & 300  & 0.3 \\
Payload (estimado)   & 800  & 3.0 \\
\midrule
\textbf{Total}       & \textbf{3000} & \textbf{5.5} \\
\bottomrule
\end{tabular}
\end{center}

Masa total $\leq 4$ kg (limite CDS para 3U). Consumo idle $\leq 6$ W
(promedio de generacion solar).

% =====================================================================
\chapter{Fases de vida del satelite}
\label{ch:fases}

NASA define oficialmente \textbf{7 fases} en el ciclo de vida de una
mision. Estas son las de aplicacion universal:

\section{Pre-A: Estudio de concepto}

\textbf{Que hace}: definir si la mision es posible y vale la pena.

\textbf{Output}: documento de concepto, presupuesto preliminar,
identificacion de tecnologia critica.

\textbf{Duracion tipica CubeSat}: 3--6 meses.

\section{Fase A: Diseño preliminar}

\textbf{Que hace}: arquitectura del sistema, decisiones top-down.

\textbf{Output}: arquitectura por subsistema, requirements de alto nivel,
diseño preliminar.

\textbf{Hito}: \textbf{SRR} (System Requirements Review) y \textbf{PDR}
(Preliminary Design Review).

\textbf{Duracion CubeSat}: 6--12 meses.

\section{Fase B: Diseño detallado}

\textbf{Que hace}: diseño detallado de cada subsistema, prototipos.

\textbf{Output}: planos, esquematicos, codigo prototipo, BOM completo.

\textbf{Hito}: \textbf{CDR} (Critical Design Review).

\textbf{Duracion}: 12 meses.

\section{Fase C: Fabricacion}

\textbf{Que hace}: construir el modelo de ingenieria (EM) y modelo de
vuelo (FM).

\textbf{Output}: hardware ensamblado, software completo.

\textbf{Hito}: \textbf{TRR} (Test Readiness Review).

\textbf{Duracion}: 6--12 meses.

\section{Fase D: Integracion y test}

\textbf{Que hace}: integrar subsistemas, hacer todo el test campaign.

\textbf{Output}: satelite testeado contra todos los estandares
ambientales.

\textbf{Hito}: \textbf{FRR} (Flight Readiness Review).

\textbf{Duracion}: 6 meses.

\section{Fase E: Operaciones}

\textbf{Que hace}: el satelite vuela y se opera desde tierra.

\textbf{Output}: datos de mision, descubrimientos cientificos.

\textbf{Duracion}: 1--5 años para CubeSat.

\section{Fase F: Disposicion}

\textbf{Que hace}: deorbit (CubeSat se quema en re-entrada) o passivacion.

\textbf{Output}: cierre de la mision.

\textbf{Duracion}: gradual, automatica por drag atmosferico.

\section{Donde estas hoy en el INTISAT}

\begin{payload}
Probable estado actual del INTISAT (a confirmar con el lider de proyecto):
\begin{itemize}[leftmargin=1em]
  \item Fase A casi cerrada (arquitectura definida).
  \item Fase B en curso (diseño detallado de subsistemas).
  \item Tu payload de microscopia esta en transicion de Fase B (diseño)
        a Fase C (fabricacion / bring-up).
\end{itemize}
\end{payload}

% =====================================================================
\chapter{Como llega al espacio}
\label{ch:lanzamiento}

\section{Rideshare --- el modelo CubeSat}

Los CubeSats NO se lanzan solos. Comparten lanzamiento con un satelite
grande (\emph{rideshare}). El lanzador principal lleva su carga primaria a
una orbita objetivo y los CubeSats se despliegan en orbitas adyacentes.

\textbf{Costo tipico}: USD 50,000 -- 200,000 para llevar un 3U a LEO.
Servicios:
\begin{itemize}[leftmargin=2em]
  \item NanoRacks (USA): el lider, desplegador desde la ISS.
  \item SpaceX Transporter: rideshare mensual a Sun-Sync.
  \item Rocket Lab Electron: dedicado a smallsat.
  \item ISRO PSLV: opcion mas barata, fiabilidad probada.
  \item Arianespace Vega: alternativa europea.
\end{itemize}

\section{Despliegue desde el lanzador}

El CubeSat va comprimido en un \textbf{deployer} (caja con resorte). El
deployer mas comun es el \textbf{P-POD} (Poly Picosatellite Orbital
Deployer):

\begin{itemize}[leftmargin=2em]
  \item Caja de aluminio que mantiene el CubeSat sujeto durante el
        lanzamiento.
  \item Cuando el lanzador llega a la orbita objetivo, el P-POD se abre.
  \item Un resorte empuja al CubeSat hacia afuera con velocidad relativa
        de $\sim 1.5$ m/s.
  \item Los rails del CubeSat deslizan en las guias del P-POD.
\end{itemize}

\section{Despliegue desde la ISS}

Tambien se puede subir el CubeSat como carga a la ISS y desplegarlo desde
ahi:
\begin{itemize}[leftmargin=2em]
  \item La carga llega a la ISS via Dragon (SpaceX), Cygnus o HTV.
  \item Astronautas lo cargan en el deployer NanoRacks dentro del modulo
        Kibo (japones).
  \item Brazo robotico extrae el deployer al exterior.
  \item Se dispara automaticamente.
\end{itemize}

\textbf{Ventaja}: el lanzamiento ya esta financiado por la NASA / ISS;
solo se paga el deploy ($\sim$ USD 50,000 / kg).

\textbf{Desventaja}: orbita fija ($\sim$ 400 km, 51.6° de inclinacion).

\section{Que pasa durante el lanzamiento}

\textbf{Cargas mecanicas}:
\begin{itemize}[leftmargin=2em]
  \item Aceleracion estatica: hasta 8 g.
  \item Vibracion aleatoria: 6--14 g RMS, 20-2000 Hz.
  \item Shock al deployment: hasta 1500 g.
\end{itemize}

\textbf{Cargas termicas}:
\begin{itemize}[leftmargin=2em]
  \item Temperatura del fairing: -100 a +120 °C.
  \item Vacio: $10^{-6}$ Torr.
\end{itemize}

Por eso el test campaign pre-vuelo es tan exhaustivo.

% =====================================================================
\chapter{Orbita LEO 101}
\label{ch:orbita}

\section{Definicion de LEO}

\textbf{LEO} = Low Earth Orbit. Altura entre 200 y 2000 km. La mayoria de
los CubeSats vuelan en LEO porque:
\begin{itemize}[leftmargin=2em]
  \item Mas barato llegar (menos energia).
  \item Buena resolucion para EO.
  \item Comms con ground station factibles.
  \item Vida util limitada (drag atmosferico genera deorbit natural).
\end{itemize}

\section{Parametros orbitales}

\begin{center}
\begin{tabular}{lll}
\toprule
\textbf{Parametro} & \textbf{Tipico} & \textbf{Que mide} \\
\midrule
Altura (h)              & 400--600 km   & distancia a la superficie \\
Periodo orbital         & 90--96 min    & tiempo de una vuelta \\
Inclinacion             & 51.6° (ISS) o 97° (SSO) & angulo respecto al ecuador \\
Eccentricity            & $\sim 0$ (circular) & forma del ovalo \\
Velocidad               & 7.7 km/s     & velocidad lineal \\
RAAN                    & 0--360°       & posicion del nodo ascendente \\
\bottomrule
\end{tabular}
\end{center}

\section{Sun-Synchronous Orbit (SSO)}

La orbita preferida para Earth Observation: el satelite pasa siempre
sobre el mismo punto a la misma hora local del dia. Garantiza
iluminacion solar constante.

Inclinacion tipica: 97°-98° (casi polar pero retrograda).

\section{Eclipses}

En LEO, el satelite entra al \textbf{cono de sombra} de la Tierra
aproximadamente \textbf{$\sim 35$\%} del tiempo (orbita iluminada $\sim 65$\%).

\begin{calculo}
Para una orbita a 500 km con periodo de 94 min:
\begin{itemize}[leftmargin=1em]
  \item Tiempo iluminado: $\sim 61$ min.
  \item Tiempo en eclipse: $\sim 33$ min.
  \item Frecuencia: 15--16 eclipses por dia.
\end{itemize}
\end{calculo}

\textbf{Implicancia}: la bateria tiene que aguantar 33 min sin generacion
solar. Eso dimensiona la capacidad de la bateria.

\section{Pases sobre la ground station}

Una ground station fija ve el satelite cuando pasa "cerca" en su orbita:
\begin{itemize}[leftmargin=2em]
  \item Pase tipico: $5-15$ minutos.
  \item Frecuencia: 4-6 pases por dia para una ground station ecuatorial,
        12+ para una polar.
  \item Ventana operativa: el tiempo en el que la antena ve al satelite
        con elevacion $> 10$° sobre el horizonte.
\end{itemize}

\section{Vida util y deorbit}

El drag atmosferico (aire residual a $400$ km) frena el satelite. Como
consecuencia:
\begin{itemize}[leftmargin=2em]
  \item Altura disminuye gradualmente.
  \item Cuando llega a $\sim 100$ km, entra en re-entrada destructiva.
  \item Tiempo total: 2--10 años desde 400--600 km segun el ciclo solar.
\end{itemize}

\textbf{Es bueno}: cumple la regla del "25 años max" para evitar basura
espacial.

% =====================================================================
\chapter{El ecosistema}
\label{ch:ecosistema}

\section{Agencias espaciales}

\begin{center}
\begin{tabular}{lp{10cm}}
\toprule
\textbf{Agencia} & \textbf{Rol en CubeSat} \\
\midrule
NASA (USA)         & Programa CSLI (CubeSat Launch Initiative): lanza
                     gratis CubeSats educativos. \\
ESA (Europa)       & Programa Fly Your Satellite + Phi-Lab para AI. \\
JAXA (Japon)       & Despliega desde la ISS via Kibo. \\
ISRO (India)       & Lanzador PSLV con CubeSats rideshare. \\
CONIDA (Peru)      & Agencia espacial nacional. Opera PeruSAT-1. \\
AEM (Mexico)       & similar. \\
CONAE (Argentina)  & Mas grande regionalmente. SAOCOM. \\
\bottomrule
\end{tabular}
\end{center}

\section{Universidades}

Las grandes universidades aeroespaciales tienen sus propios CubeSats:

\begin{itemize}[leftmargin=2em]
  \item Stanford, MIT, Cal Poly (USA).
  \item TU Delft, Aalborg, ETH Zurich (Europa).
  \item NTNU (Noruega) --- HYPSO.
  \item Tokyo Tech (Japon).
  \item KAIST (Corea).
  \item UNI, PUCP (Peru).
  \item UNLP, ITBA (Argentina).
\end{itemize}

\section{Empresas comerciales}

\textbf{NewSpace startups} que cambiaron el juego:

\begin{center}
\begin{tabular}{lp{10cm}}
\toprule
\textbf{Empresa} & \textbf{Que hace} \\
\midrule
Planet Labs (USA)        & 200+ Doves: imagenes diarias de toda la Tierra. \\
Spire Global (USA)       & Constelacion para meteorologia y AIS. \\
Satellogic (Argentina)   & EO de alta resolucion. \\
ICEYE (Finlandia)        & SAR (radar) en CubeSats. \\
Innova Space (Argentina) & IoT desde el espacio. \\
GomSpace (Dinamarca)     & Componentes y satellites para terceros. \\
ISIS (Holanda)           & Lider en componentes de CubeSat. \\
EnduroSat (Bulgaria)     & Componentes y misiones. \\
AAC Clyde Space (UK)     & Servicios end-to-end. \\
KP Labs (Polonia)        & AI on-board (Intuition-1). \\
Aiko Space (Italia)      & Software AI para misiones. \\
\bottomrule
\end{tabular}
\end{center}

\section{Lanzadores}

\begin{center}
\begin{tabular}{lp{8cm}p{2cm}}
\toprule
\textbf{Lanzador} & \textbf{Servicio} & \textbf{Costo} \\
\midrule
SpaceX Falcon 9 (Transporter) & rideshare mensual a SSO          & $\sim$ 6 k\$/kg \\
Rocket Lab Electron            & smallsat dedicado                & $\sim$ 25 k\$/kg \\
ISRO PSLV                      & rideshare desde India            & $\sim$ 10 k\$/kg \\
Arianespace Vega-C             & rideshare europeo                & $\sim$ 15 k\$/kg \\
NASA CSLI                       & GRATIS para misiones educativas  & 0 \\
NanoRacks (desde ISS)           & deploy desde la ISS              & $\sim$ 50 k\$/kg \\
\bottomrule
\end{tabular}
\end{center}

\section{Coordinacion regulatoria}

\begin{itemize}[leftmargin=2em]
  \item \textbf{ITU} (UIT): coordina espectro radioelectrico a nivel
        internacional.
  \item \textbf{IARU}: coordinacion para CubeSats en banda amateur (UHF).
  \item Autoridades nacionales: ANATEL (Brasil), MITC (Peru), AGESIC
        (Uruguay), etc.
  \item \textbf{ITU-R} registrar frecuencias.
\end{itemize}

% =====================================================================
\chapter{Estandares clave}
\label{ch:estandares}

\section{CalPoly CubeSat Design Specification (CDS)}

\textbf{El} estandar fundacional. Define:
\begin{itemize}[leftmargin=2em]
  \item Dimensiones mecanicas (1U, 3U, 6U, etc.).
  \item Masa maxima.
  \item Materiales aceptables.
  \item Interfaz con el deployer.
  \item Requisitos de testing pre-deployment.
\end{itemize}

URL: \url{https://www.cubesat.org/cds-announcement}

Version actual: \textbf{Rev 14} (2022).

\section{ECSS --- European Cooperation for Space Standardization}

Conjunto de estandares europeos. Los relevantes para CubeSat:
\begin{itemize}[leftmargin=2em]
  \item ECSS-E-ST-10-03C --- Verification \& Testing.
  \item ECSS-Q-ST-30C --- Dependability.
  \item ECSS-Q-ST-60C --- Electrical, Electronic, Electromechanical.
  \item ECSS-E-ST-70-11C --- Operability.
\end{itemize}

\section{NASA GEVS --- GSFC-STD-7000}

General Environmental Verification Standard de Goddard Space Flight
Center. Define tests ambientales:
\begin{itemize}[leftmargin=2em]
  \item Vibracion (sine sweep, random, shock).
  \item Termico vacuum (TVAC).
  \item EMC/EMI.
  \item Magnetic.
\end{itemize}

\section{MIL-STD-810}

Estandar militar de tests ambientales. Mas estricto que GEVS en muchos
aspectos. Usado por CubeSats militares.

\section{CCSDS --- Consultative Committee for Space Data Systems}

Estandares de comunicacion y datos espaciales:
\begin{itemize}[leftmargin=2em]
  \item CCSDS Space Packet: empaquetado de telemetria.
  \item CCSDS File Delivery Protocol (CFDP): transferencia de archivos.
  \item CCSDS Encapsulation Service.
\end{itemize}

% =====================================================================
\chapter{Vocabulario --- glosario completo}
\label{ch:glosario}

Lista de las siglas y conceptos que vas a encontrar todo el tiempo:

\section*{A}

\begin{itemize}[leftmargin=2em]
  \item \textbf{ADCS}: Attitude Determination and Control System. El que
        controla la orientacion del satelite.
  \item \textbf{AIT}: Assembly, Integration and Test. La fase de
        construir, integrar y probar.
  \item \textbf{AGN}: Anomaly Generation Note. Reporte de anomalia operativa.
  \item \textbf{AOCS}: Attitude and Orbit Control System (sinonimo de ADCS).
  \item \textbf{AOS / LOS}: Acquisition of Signal / Loss of Signal. Inicio
        y fin de un pase.
\end{itemize}

\section*{B--C}

\begin{itemize}[leftmargin=2em]
  \item \textbf{BOL / EOL}: Beginning / End of Life. Inicio y fin de
        mision.
  \item \textbf{CAN bus}: Controller Area Network. Bus serial robusto.
  \item \textbf{cFS}: Core Flight System. Framework de flight software de NASA.
  \item \textbf{CCSDS}: Consultative Committee for Space Data Systems.
  \item \textbf{CDR}: Critical Design Review.
  \item \textbf{CDS}: CubeSat Design Specification.
  \item \textbf{CG}: Centro de Gravedad.
  \item \textbf{COTS}: Commercial Off-The-Shelf. Componentes comerciales no
        certificados para vuelo.
  \item \textbf{CSP}: CubeSat Space Protocol. Stack de comunicacion de
        GomSpace.
  \item \textbf{CSLI}: CubeSat Launch Initiative (NASA).
\end{itemize}

\section*{D--E}

\begin{itemize}[leftmargin=2em]
  \item \textbf{DCS / SCADA}: estandares industriales (no aeroespaciales,
        pero similares).
  \item \textbf{DoD}: Depth of Discharge de bateria.
  \item \textbf{DPC}: Differential Phase Contrast (microscopia).
  \item \textbf{Downlink}: comunicacion del satelite hacia la tierra.
  \item \textbf{Deorbit}: re-entrada controlada o natural a la atmosfera.
  \item \textbf{EPS}: Electrical Power System.
\end{itemize}

\section*{F--G}

\begin{itemize}[leftmargin=2em]
  \item \textbf{FCC}: Federal Communications Commission (USA).
  \item \textbf{FDIR}: Fault Detection, Isolation, Recovery.
  \item \textbf{FOV}: Field of View.
  \item \textbf{FPM}: Fourier Ptychographic Microscopy.
  \item \textbf{FPR}: Flight Readiness Review.
  \item \textbf{FSW}: Flight Software.
  \item \textbf{GEO}: Geostationary Earth Orbit (35 786 km).
  \item \textbf{GEVS}: General Environmental Verification Standard.
  \item \textbf{GMAT}: General Mission Analysis Tool (NASA).
  \item \textbf{GPS}: Global Positioning System.
  \item \textbf{Ground segment / station}: la parte terrestre del sistema.
\end{itemize}

\section*{H--L}

\begin{itemize}[leftmargin=2em]
  \item \textbf{HK}: Housekeeping. Telemetria de estado.
  \item \textbf{I2C / SPI / UART}: protocolos seriales.
  \item \textbf{IARU}: International Amateur Radio Union.
  \item \textbf{ITU}: International Telecommunication Union.
  \item \textbf{LEO}: Low Earth Orbit (200-2000 km).
  \item \textbf{LDO}: Low Dropout regulator.
  \item \textbf{LoS}: Line of Sight.
\end{itemize}

\section*{M--O}

\begin{itemize}[leftmargin=2em]
  \item \textbf{MEO}: Medium Earth Orbit (2000-35 786 km).
  \item \textbf{MIL-STD-810}: estandar militar de tests ambientales.
  \item \textbf{MPPT}: Maximum Power Point Tracker (para paneles solares).
  \item \textbf{NA}: Numerical Aperture (optica).
  \item \textbf{OBC}: On-Board Computer.
  \item \textbf{OBDH}: On-Board Data Handling.
  \item \textbf{Orbit}: orbita.
  \item \textbf{OTA}: Over The Air (update).
\end{itemize}

\section*{P--R}

\begin{itemize}[leftmargin=2em]
  \item \textbf{Pass}: pase del satelite sobre la ground station.
  \item \textbf{Payload}: la carga util cientifica/comercial.
  \item \textbf{PC-104}: estandar de bus electrico/mecanico para CubeSat.
  \item \textbf{PDR}: Preliminary Design Review.
  \item \textbf{PMIC}: Power Management Integrated Circuit.
  \item \textbf{P-POD}: Poly Picosatellite Orbital Deployer.
  \item \textbf{PSP}: Platform Support Package (en cFS).
  \item \textbf{RAAN}: Right Ascension of Ascending Node.
  \item \textbf{RF}: Radio Frequency.
\end{itemize}

\section*{S--Z}

\begin{itemize}[leftmargin=2em]
  \item \textbf{SAA}: South Atlantic Anomaly. Zona de alta radiacion sobre
        el Atlantico Sur. Los satellites suelen entrar a SAFE\_MODE.
  \item \textbf{SAR}: Synthetic Aperture Radar.
  \item \textbf{S-Band}: 2.2-2.4 GHz para comms.
  \item \textbf{SEU / TID}: Single Event Upset / Total Ionizing Dose
        (efectos de radiacion).
  \item \textbf{SoC}: System on Chip.
  \item \textbf{SSO}: Sun-Synchronous Orbit.
  \item \textbf{Telemetry}: datos del estado del satelite.
  \item \textbf{TLE}: Two-Line Element. Formato estandar de elementos
        orbitales.
  \item \textbf{TRL}: Technology Readiness Level (NASA 1-9).
  \item \textbf{TVAC}: Thermal Vacuum.
  \item \textbf{UHF}: Ultra High Frequency (300 MHz -- 3 GHz).
  \item \textbf{Uplink}: comunicacion de la tierra hacia el satelite.
  \item \textbf{V\&V}: Verification \& Validation.
  \item \textbf{VHF}: Very High Frequency (30-300 MHz).
  \item \textbf{X-Band}: 8-12 GHz para comms de alto throughput.
\end{itemize}

% =====================================================================
\appendix
\chapter{Apendice --- Misiones de referencia}

\section{PhiSat-1 (ESA, 2020)}

Primer satelite con IA on-board. Detecta nubes con CNN en Intel Movidius
Myriad 2 antes de bajar las imagenes. Tu INTISAT esta en la misma linea
metodologica.

\textbf{Lecciones aprendidas}:
\begin{itemize}[leftmargin=2em]
  \item Filtrar antes de bajar ahorra bitrate.
  \item Modelo pre-entrenado en GPUs grandes y desplegado en hardware
        embebido funciona.
  \item Validacion previa en TVAC es critica.
\end{itemize}

\section{HYPSO-1 (NTNU Noruega, 2022)}

CubeSat 6U con hiperspectral y RPi CM4 + IA on-board. \textbf{El hardware mas
cercano al INTISAT en mision real}.

\textbf{Lecciones}:
\begin{itemize}[leftmargin=2em]
  \item RPi CM4 sobrevive bien en LEO.
  \item eMMC es mas confiable que SD para vuelo.
  \item El downlink es el cuello de botella, no el procesado.
\end{itemize}

\section{GeneSat-1 (NASA Ames, 2006)}

Primer CubeSat de biologia. Levo E. coli para experimentos biologicos en
microgravedad. Demuestra que microscopia biologica en CubeSat es viable.

\section{PerúSAT-1 (CONIDA, 2016)}

Satelite peruano de observacion de la Tierra. No es CubeSat (430 kg) pero
es el referente nacional. Construido por Airbus, operado por CONIDA.

\chapter{Apendice --- Recursos para profundizar}

\section{Libros}

\begin{itemize}[leftmargin=2em]
  \item \emph{Space Mission Analysis and Design} (Wertz \& Larson) ---
        el clasico. USD 200.
  \item \emph{CubeSat Engineering Notebook} (Aalborg Univ). Gratis online.
  \item \emph{CubeSat Handbook} (Cappelletti et al). Academic Press.
\end{itemize}

\section{Cursos online (gratis)}

\begin{itemize}[leftmargin=2em]
  \item \textbf{Coursera} -- "Spacecraft Dynamics and Control"
        (University of Colorado).
  \item \textbf{edX} -- "Small Satellite Engineering" (TU Delft).
  \item \textbf{NASA CubeSat 101} -- introduccion gratis en NASA.gov.
  \item \textbf{Cal Poly CubeSat Workshop} -- videos en YouTube.
\end{itemize}

\section{Comunidades}

\begin{itemize}[leftmargin=2em]
  \item Cal Poly CubeSat Workshop (anual).
  \item Small Satellite Conference USU (anual, todas las papers gratis).
  \item ESA Fly Your Satellite Programme.
  \item r/CubeSat en Reddit.
  \item Discord servers academicos.
\end{itemize}

\chapter{Apendice --- Hoja de ruta de aprendizaje}

\section{Plan de 3 meses para un junior en CubeSat}

\textbf{Mes 1}: este manual M0 completo + Manual de Electronica Practica.

\textbf{Mes 2}: Manual M2 (EPS) + Manual M5 (Flight Software).

\textbf{Mes 3}: Manual M7 (Testing) + especializacion en el subsistema de
interes.

\section{Plan de 1 año para construir un CubeSat completo}

\begin{itemize}[leftmargin=2em]
  \item \textbf{T+0 a T+3 meses}: Fase A. Concepto, requirements.
  \item \textbf{T+3 a T+8 meses}: Fase B. Diseño detallado, prototipos.
  \item \textbf{T+8 a T+11 meses}: Fase C-D. Fabricacion e integracion.
  \item \textbf{T+11 a T+12 meses}: Testing y FRR.
  \item \textbf{T+12}: lanzamiento.
\end{itemize}

\textbf{En la practica}: para un equipo academico de 8-10 personas, son
2-3 años reales. Los CubeSats academicos siempre subestiman tiempo.

\backmatter

\chapter*{Bibliografia}
\addcontentsline{toc}{chapter}{Bibliografia}

\begin{enumerate}[leftmargin=2em]
\item Cal Poly. \emph{CubeSat Design Specification Rev 14}.
  \url{https://www.cubesat.org}, 2022.
\item Wertz J, Everett D, Puschell J. \emph{Space Mission Engineering: The
  New SMAD}. Microcosm Press, 2011.
\item Bouwmeester J, Guo J. State of the art of CubeSats. \emph{Acta
  Astronautica}, 67(7-8), 854-862, 2010.
\item Swartwout M. The first one hundred CubeSats. \emph{Journal of Small
  Satellites}, 2013.
\item NASA. \emph{CubeSat 101 - Basic Concepts and Processes for
  First-Time CubeSat Developers}. NASA Headquarters, 2017.
\item ESA. \emph{Fly Your Satellite! Programme Documentation}.
  \url{https://www.esa.int/Education/CubeSats}, 2022.
\item Furano G, et al. AI in Space: Opportunities and Challenges.
  \emph{ESA Advanced Concepts Team Report}, 2020.
\item Marcuccio S, et al. Lessons from PhiSat-1. \emph{IEEE Aerospace
  Conference}, 2021.
\end{enumerate}

\end{document}
